\documentclass[10pt,a4paper]{amsart}
\usepackage[margin=19mm]{geometry}
\usepackage{amsmath,amssymb,mathtools}
\usepackage[scr=rsfs,cal=euler]{mathalfa}
\usepackage{xfrac}
\usepackage[unicode,hidelinks,bookmarks=false]{hyperref}
\newcommand{\ie}{i.e.\ }
\newcommand{\cf}{cf.\ }
\newcommand{\resp}{respectively\ }
\newcommand{\Subsection}{\subsection}
\providecommand{\nobreakdash}{\nobreak\hbox{-}}
\setlength{\emergencystretch}{2em}
\multlinegap0pt
\usepackage{float}
\usepackage{mathtools}
\usepackage{amssymb}
\usepackage{algorithm}\usepackage{algpseudocode}
\usepackage{siunitx}
\newtheorem{assumption}{Assumption}
\newtheorem{lemma}{Lemma}
\newcommand{\bR}{\mathbb{R}}  
\title{A novel sequential method for building upper and lower bounds of moments of distributions}
\author{Solal Martin, Emilie Chouzenoux, Victor Elvira}
\date{Source: arXiv:2512.01761v2 (CC BY 4.0)}
\begin{document}
\maketitle

\noindent\emph{Source and licence.} A novel sequential method for building upper and lower bounds of moments of distributions. Version arXiv:2512.01761v2 (CC BY 4.0), distributed by arXiv under the Creative Commons Attribution 4.0 International licence, http://creativecommons.org/licenses/by/4.0/, as recorded in the arXiv licence metadata for this exact version and shown on the abstract page https://arxiv.org/abs/2512.01761v2. Full licence notice: \texttt{licenses/arxiv-2512.01761-CC-BY-4.0.txt}.\par\smallskip

\noindent\textit{Editorial scope.} Counted unit: the article's lemma lem:bounds\_p2\_over\_q (source0.tex lines 3356--3407). The article's complete original proof of it is delivered with it (source0.tex lines 3409--3607, one \texttt{proof} environment of 199 lines). No mathematics is written, repaired or re-derived: the statement and the proof are the article's own lines, in the article's own notation, and no retained line is substituted or re-flowed.  Retained uncounted as the source-local context the proof needs: lines 226--229; lines 231--234; lines 246--253; lines 284--300; lines 326--332; lines 354--371; lines 3332--3336.  Nothing was omitted from any retained span, so the package carries no omission record.  No retained line contains a citation, so the wrapper carries no bibliography and no bibliography entry.  No retained line contains a figure environment, an image-inclusion command or drawing code, so no image is delivered and the package declares no asset dependency.  Editorial formatting only, in addition: an AMS \texttt{amsart} preamble that restates the article's own declarations and macros used by the retained lines, namely the environments \texttt{assumption}, \texttt{lemma} on separate counters, together with the standard packages those lines need.  The retained environments print in this document as Assumption 1, Lemma 1 in order of appearance, whereas the article prints the same items with the numbering of its own full document.  The complete native source of the article as distributed in the arXiv e-print for this exact version is bundled unchanged as \texttt{sources/190150/source0.tex}.
\begin{equation}
(\forall x \in \bR) \quad \pi(x) \geq \underline{b}(x;t) \quad \text{and} \quad \pi(t) = \underline{b}(t;t).
\label{eq_minor}
\end{equation}

\begin{equation}
(\forall x \in \bR) \quad \pi(x) \leq \overline{b}(x;t) \quad \text{and} \quad \pi(t) = \overline{b}(t;t).
\label{eq_major}
\end{equation}

\begin{assumption}
\label{assump1}
For every $t \in \mathbb{R}$, there exists $\beta(t) > 0$ such that 
\begin{equation}
(\forall x \in \mathbb{R}) \quad \phi(x) \leq \phi(t) + \dot{\phi}(t) (x-t) + \frac{\beta(t)}{2} (x-t)^2,
\end{equation}
where $\dot{\phi}$ denotes the first order derivative of $\phi$. 
\end{assumption}

\begin{lemma}
\label{lem:ass1}
Under Assumption~\ref{assump1}, the minoration conditions \eqref{eq_minor} hold with
\begin{equation}
(\forall (x,t) \in \bR^2) \quad \underline{b}(x;t) = 
\underline{C}(t) g(x; \underline{u}(t), \underline{\sigma}(t))
\end{equation}
using the following variance, mean and scale parameters:
\begin{equation}
(\forall t \in \bR) \quad
\begin{cases}
\underline{\sigma}(t) & = 1/\sqrt{\beta(t)}, \\
\underline{u}(t) & = t -\underline{\sigma}(t)^2 \dot{\phi}(t), \\
\underline{C}(t) & = \sqrt{2 \pi \underline{\sigma}(t)^2} \exp\left(- \phi(t) + (\underline{\sigma}(t)^2/2) (\dot{\phi}(t))^2\right). \\
\end{cases}
\end{equation}
\end{lemma}

\begin{assumption}
\label{assump2}
For every $t \in \mathbb{R}$, there exists $\nu(t)>0$ such that
\begin{equation}
(\forall x \in \mathbb{R}) \quad \phi(x) \geq \phi(t) + \dot{\phi}(t) (x-t) + \frac{\nu(t)}{2} (x-t)^2
\end{equation}
\end{assumption}

\begin{lemma}
\label{lem:ass2}
Under Assumption~\ref{assump2}, the majoration conditions \eqref{eq_major} are satisfied, as soon as
\begin{equation}
(\forall (x,t) \in \bR^2) \quad
\overline{b}(x;t)    = \overline{C}(t) g(x; \overline{u}(t), \overline{\sigma}(t)),
\end{equation}
with variance, mean and scale parameters given by
\begin{equation}
\label{eq:uppermean_variance_constant_gauss}
(\forall t \in \bR) \quad
\begin{cases}
\overline{\sigma}(t) & = 1/\sqrt{\nu(t)},\\
\overline{u}(t) & = t -\overline{\sigma}(t)^2 \dot{\phi}(t), \\
\overline{C}(t) & = \sqrt{2 \pi \overline{\sigma}(t)^2} \exp\left(- \phi(t) + (\overline{\sigma}(t)^2/2) (\dot{\phi}(t))^2\right). \\
\end{cases}
\end{equation}
\end{lemma}

\begin{equation}
\label{bounds_over_p}
\underline{C}(t)\, g(x;\underline{u}(t),\underline{\sigma}(t)) \leq p(x) 
\leq \overline{C}(t)\, g(x;\overline{u}(t),\overline{\sigma}(t))
\end{equation}

\begin{lemma}
\label{lem:bounds_p2_over_q}
Assume that, for all $t \in \mathbb{R}$, $(\underline{\sigma}(t)^2,\overline{\sigma}(t)^2)\in (0,2\theta^2)$, then, for every $(x,t) \in \mathbb{R}^2$,
\begin{equation}
\underline{D}(t)\, g(x;\underline{v}(t),\underline{\rho}(t)) \leq
\frac{p(x)^2}{q(x)}\leq \overline{D}(t)\, g(x;\overline{v}(t),\overline{\rho}(t)),
\label{eq:boundp2q}
\end{equation}
where
\begin{equation}
\begin{cases}
\underline{\rho}(t)^2
=
\left(\dfrac{2}{\underline{\sigma}(t)^2}-\dfrac{1}{\theta^2}\right)^{-1},\\[2ex]
\underline{v}(t)
=
\underline{\rho}(t)^2
\left(
\dfrac{2\underline{u}(t)}{\underline{\sigma}(t)^2}
-\dfrac{\mu}{\theta^2}
\right),\\[2ex]
\underline{D}(t)
=
\underline{C}(t)^2
\dfrac{\underline{\rho}(t)}{\underline{\sigma}(t)^2}\theta
\exp\left(
-\dfrac{\underline{u}(t)^2}{\underline{\sigma}(t)^2}
+\dfrac{\mu^2}{2\theta^2}
+\dfrac{\underline{v}(t)^2}{2\underline{\rho}(t)^2}
\right),\\
\overline{\rho}(t)^2
=
\left(\dfrac{2}{\overline{\sigma}(t)^2}-\dfrac{1}{\theta^2}\right)^{-1},\\[2ex]
\overline{v}(t)
=
\overline{\rho}(t)^2
\left(
\dfrac{2\overline{u}(t)}{\overline{\sigma}(t)^2}
-\dfrac{\mu}{\theta^2}
\right),\\[2ex]
\overline{D}(t)
=
\overline{C}(t)^2
\dfrac{\overline{\rho}(t)}{\overline{\sigma}(t)^2}\theta
\exp\left(
-\dfrac{\overline{u}(t)^2}{\overline{\sigma}(t)^2}
+\dfrac{\mu^2}{2\theta^2}
+\dfrac{\overline{v}(t)^2}{2\overline{\rho}(t)^2}
\right).
\end{cases}
\end{equation}
\end{lemma}

\begin{proof}
We have
$$
(\forall x \in \mathbb{R}) \quad q(x)=\frac{1}{\sqrt{2\pi \theta^2}}
\exp\left(-\frac{(x-\mu)^2}{2\theta^2}\right),
$$
hence
\begin{equation}
(\forall x \in \mathbb{R}) \quad
\frac{p(x)^2}{q(x)}
=
\exp(-2\phi(x))\sqrt{2\pi \theta^2}
\exp\left(\frac{(x-\mu)^2}{2\theta^2}\right).
\end{equation}
Using the left inequality in~\eqref{bounds_over_p}, and the above expression of $q$ leads, for every $t \in \mathbb{R}$, to
% Under Assumption~\ref{assump1}, Lemma~\ref{lem:ass1} yields
% \begin{equation}
% \phi(x)
% \leq
% \frac{1}{2\underline{\sigma}(t)^2}(x-\underline{u}(t))^2
% +\phi(t)
% -\frac{\underline{\sigma}(t)^2}{2}(\dot{\phi}(t))^2.
% \end{equation}
% Multiplying by $-2$, we get
% \begin{equation}
% -2\phi(x)
% \geq
% -\frac{1}{\underline{\sigma}(t)^2}(x-\underline{u}(t))^2
% -2\phi(t)
% +\underline{\sigma}(t)^2(\dot{\phi}(t))^2.
% \end{equation}
% Hence
\begin{align}
(\forall x \in \mathbb{R}) \quad \frac{p(x)^2}{q(x)}
&\geq
\sqrt{2\pi \theta^2}
\exp\left(
-\frac{1}{\underline{\sigma}(t)^2}(x-\underline{u}(t))^2
+\frac{(x-\mu)^2}{2\theta^2}
-2\phi(t)
+\underline{\sigma}(t)^2(\dot{\phi}(t))^2
\right).
\end{align}
Expanding the quadratic terms gives, for every $(x,t) \in \mathbb{R}^2$,
\begin{multline}
-\frac{1}{\underline{\sigma}(t)^2}(x-\underline{u}(t))^2
+\frac{(x-\mu)^2}{2\theta^2}
=
-\left(
\frac{1}{\underline{\sigma}(t)^2}
-\frac{1}{2\theta^2}
\right)x^2
+
\left(
\frac{2\underline{u}(t)}{\underline{\sigma}(t)^2}
-\frac{\mu}{\theta^2}
\right)x  
-\frac{\underline{u}(t)^2}{\underline{\sigma}(t)^2}
+\frac{\mu^2}{2\theta^2}.
\label{eq_170}
\end{multline}
Since we assumed that, for every $t \in \mathbb{R}$, $\underline{\sigma}(t)^2<2\theta^2$, then
$$
\frac{1}{\underline{\sigma}(t)^2}-\frac{1}{2\theta^2}>0.
$$
Therefore, by completing the square, we can rewrite \eqref{eq_170} for every $(x,t) \in \mathbb{R}^2$ as
\begin{equation}
-\frac{1}{\underline{\sigma}(t)^2}(x-\underline{u}(t))^2
+\frac{(x-\mu)^2}{2\theta^2}
=
-\frac{1}{2\underline{\rho}(t)^2}(x-\underline{v}(t))^2
-\frac{\underline{u}(t)^2}{\underline{\sigma}(t)^2}
+\frac{\mu^2}{2\theta^2}
+\frac{\underline{v}(t)^2}{2\underline{\rho}(t)^2},
\end{equation}
with the notation
\begin{equation}
(\forall t \in \mathbb{R}) \quad \underline{\rho}(t)^2
=
\left(\frac{2}{\underline{\sigma}(t)^2}-\frac{1}{\theta^2}\right)^{-1}
\end{equation}
and
\begin{equation}
(\forall t \in \mathbb{R}) \quad \underline{v}(t)
=
\underline{\rho}(t)^2
\left(
\frac{2\underline{u}(t)}{\underline{\sigma}(t)^2}
-\frac{\mu}{\theta^2}
\right).
\end{equation}
Thus, for every $(x,t) \in \mathbb{R}^2$,
\begin{align}
\frac{p(x)^2}{q(x)}
&\geq
\sqrt{2\pi \theta^2}
\exp\left(
-2\phi(t)
+\underline{\sigma}(t)^2(\dot{\phi}(t))^2
-\frac{\underline{u}(t)^2}{\underline{\sigma}(t)^2}
+\frac{\mu^2}{2\theta^2}
+\frac{\underline{v}(t)^2}{2\underline{\rho}(t)^2}
\right)
\exp\left(
-\frac{(x-\underline{v}(t))^2}{2\underline{\rho}(t)^2}
\right).
\end{align}
    Using the expression of $g(x;\underline{v}(t),\underline{\rho}(t))$ proves the left inequality in \eqref{eq:boundp2q}. The right inequality in \eqref{eq:boundp2q} is proved in a similar manner, using the right inequality in \eqref{bounds_over_p} and the assumption on the range of $\overline{\sigma}(t)$.

% $$
% g(x;\underline{v}(t),\underline{\rho}(t))
% =
% \frac{1}{\sqrt{2\pi}\underline{\rho}(t)}
% \exp\left(
% -\frac{(x-\underline{v}(t))^2}{2\underline{\rho}(t)^2}
% \right),
% $$
% we obtain
% \begin{equation}
% \frac{p(x)^2}{q(x)}
% \geq
% \underline{D}(t)\,g(x;\underline{v}(t),\underline{\rho}(t)),
% \end{equation}
% with
% \begin{equation}
% \underline{D}(t)
% =
% 2\pi\sigma\underline{\rho}(t)
% \exp\left(
% -2\phi(t)
% +\underline{\sigma}(t)^2(\dot{\phi}(t))^2
% -\frac{\underline{u}(t)^2}{\underline{\sigma}(t)^2}
% +\frac{\mu^2}{2\sigma^2}
% +\frac{\underline{v}(t)^2}{2\underline{\rho}(t)^2}
% \right).
% \end{equation}

% Since
% \begin{equation}
% \underline{C}(t)
% =
% \sqrt{2\pi \underline{\sigma}(t)^2}
% \exp\left(
% -\phi(t)+\frac{\underline{\sigma}(t)^2}{2}(\dot{\phi}(t))^2
% \right),
% \end{equation}
% this can also be rewritten as
% \begin{equation}
% \underline{D}(t)
% =
% \underline{C}(t)^2
% \dfrac{\underline{\rho}(t)}{\underline{\sigma}(t)^2}\sigma
% \exp\left(
% -\dfrac{\underline{u}(t)^2}{\underline{\sigma}(t)^2}
% +\dfrac{\mu^2}{2\sigma^2}
% +\dfrac{\underline{v}(t)^2}{2\underline{\rho}(t)^2}
% \right).
% \end{equation}

% The proof of the upper bound follows exactly the same lines, using Assumption~\ref{assump2} and Lemma~\ref{lem:ass2}. In this case,
% \begin{equation}
% \phi(x)
% \geq
% \frac{1}{2\overline{\sigma}(t)^2}(x-\overline{u}(t))^2
% +\phi(t)
% -\frac{\overline{\sigma}(t)^2}{2}(\dot{\phi}(t))^2,
% \end{equation}
% which implies
% \begin{equation}
% \frac{p(x)^2}{q(x)}
% \leq
% \overline{D}(t)\,g(x;\overline{v}(t),\overline{\rho}(t)),
% \end{equation}
% with
% \begin{equation}
% \begin{cases}
% \overline{\rho}(t)^2
% =
% \left(\dfrac{2}{\overline{\sigma}(t)^2}-\dfrac{1}{\sigma^2}\right)^{-1},\\[2ex]
% \overline{v}(t)
% =
% \overline{\rho}(t)^2
% \left(
% \dfrac{2\overline{u}(t)}{\overline{\sigma}(t)^2}
% -\dfrac{\mu}{\sigma^2}
% \right),\\[2ex]
% \overline{D}(t)
% =
% \overline{C}(t)^2
% \dfrac{\overline{\rho}(t)}{\overline{\sigma}(t)^2}\sigma
% \exp\left(
% -\dfrac{\overline{u}(t)^2}{\overline{\sigma}(t)^2}
% +\dfrac{\mu^2}{2\sigma^2}
% +\dfrac{\overline{v}(t)^2}{2\overline{\rho}(t)^2}
% \right).
% \end{cases}
% \end{equation}

\end{proof}

\end{document}
